\section{Lean formalization}

\subsection{Built on the Equational Theories Project}

The Lean framework we use for the formalization is largely based on tools developed in the Equational Theories Project (ETP\@).
For clarity, we isolate code from that project into the \texttt{equational\_theories} directory, which we treat as a library on the same footing as \texttt{mathlib}, while new code goes in a separate directory \texttt{associative\_theories} (see the lakefile in \autoref{fig:lakefile}).

\begin{figure}
\lstinputlisting[basicstyle=\ttfamily]{../lakefile.toml}
\caption{\label{fig:lakefile}Configuration file \texttt{lakefile.toml}}
\end{figure}

The Lean formalization splits into several tasks.
\begin{itemize}
\item Prove that all equations of order $\leq 4$ are in a set of \num{653} equations.  This follows the \texttt{laws\_complete} theorem from the ETP\@.

\item Prove that these \num{653} equations are all equivalent to \num{122}.  This is done by establishing pairs of implications from each equation to its (lowest-numbered) representative.

\item .....

\item To disprove implications, we use finite countermodels (cf.\ the \texttt{Countermodels} directory), based on ETP tools to manipulate finite operation tables.  The structure of the files differs slightly from the ETP: given the centralized nature of this project, the countermodels are numbered sequentially instead of using their operation table as a name of the model and of the corresponding Lean theorem.  In this way the operation table is only specified once in the file.
\end{itemize}

\begin{figure}
\lstinputlisting[language=lean]{../associative_theories/Countermodels/Model10.lean}
\caption{\label{fig:lakefile}An example countermodel file}
\end{figure}

\subsection{Vampire for positive proofs}

The implications checked by ATPs must also be formalized in Lean.  For this we rely on ETP tools, which we now describe.

\subsection{End-to-end theorem}

Discuss here also how complete the verification is, what are the loose ends such as if we had to write the list of equations by hand, or similar.
